\documentclass[10pt,a4paper]{amsart}
\usepackage[margin=19mm]{geometry}
\usepackage{amsmath,amssymb,mathtools}
\usepackage[scr=rsfs,cal=euler]{mathalfa}
\usepackage{xfrac}
\usepackage[unicode,hidelinks,bookmarks=false]{hyperref}
\newcommand{\ie}{i.e.\ }
\newcommand{\cf}{cf.\ }
\newcommand{\resp}{respectively\ }
\newcommand{\Subsection}{\subsection}
\providecommand{\nobreakdash}{\nobreak\hbox{-}}
\setlength{\emergencystretch}{2em}
\multlinegap0pt
\usepackage{bm}
\usepackage{bbm}
\usepackage{dsfont}
\usepackage{xspace}
\usepackage{cancel}
\usepackage{mathtools}
\usepackage[shortlabels]{enumitem}
\usepackage{amsthm}
\makeatletter
\@ifundefined{theorem}{\newtheorem{theorem}{Theorem}}{}
\@ifundefined{definition}{\newtheorem{definition}{Definition}}{}
\@ifundefined{lemma}{\newtheorem{lemma}[theorem]{Lemma}}{}
\@ifundefined{proposition}{\newtheorem{proposition}[theorem]{Proposition}}{}
\@ifundefined{remark}{\newtheorem{remark}[theorem]{Remark}}{}
\makeatother
\newcommand\Bern[3]{B\left(#3 L^n \h{#1}{} \v{#1}{#2}\right)}
\newcommand\Fluxnum[1]{\mathcal{F}_{#1 \frac{1}{2}}^n}
\newcommand\R{\mathbb{R}}
\newcommand\dt{\Delta t}
\newcommand\h[1]{h_{#1 \frac{1}{2}}}
\renewcommand\H{\mathcal{H}_{\phi}}
\renewcommand\v[2]{v_{#1 \frac{1}{2}}^{#2}}
\title[Finite Volumes for a dissipative free boundary problem]{Finite Volumes for a dissipative free boundary problem}
\author{Cl\'ement Canc\`es, Claire Chainais-Hillairet, Am\'elie Dupouy}
\date{Source: arXiv:2509.14908v1 (CC BY 4.0)}
\begin{document}
\maketitle

\noindent\emph{Source and licence.} Finite Volumes for a dissipative free boundary problem. Version arXiv:2509.14908v1 (CC BY 4.0), distributed under the Creative Commons Attribution 4.0 International licence, as recorded in the arXiv licence metadata for this exact version and shown on the abstract page https://arxiv.org/abs/2509.14908v1. Full rights notice: licenses/arxiv-2509.14908-CC-BY-4.0.txt.\par\smallskip

\noindent\textit{Editorial scope.} Counted unit: the Proposition whose source label is \texttt{prop:disc\_energ\_dec} in the article (source0.tex lines 535--541, source label \texttt{prop:disc\_energ\_dec}), delivered together with the article's own complete original proof of it (source0.tex lines 542--606, one \texttt{proof} environment of 65 lines). No mathematics is written, repaired or re-derived: statement and proof are the article's own text line for line. Required context retained uncounted: eq:scheme (lines 297--311); eq:left\_cond\_flux (lines 298--310); eq:right\_cond\_flux (lines 298--310); eq:scheme\_u (lines 298--310); eq:def\_v (lines 319--322); lem:Bern (lines 325--335); eq:Bern (lines 331--333); rk:pi\_croiss (lines 454--456); eq:def\_htot\_disc (lines 522--532). Every definition, display and label of those lines is retained unchanged. Omitted from this package and not redistributed: 1--296, 311--318, 323--324, 334--453, 457--521, 533--534, 607--1274. Nothing inside a retained excerpt is omitted or altered: every retained body is delivered exactly as the article's lines are written, so no omission receipt of any kind accompanies this package. Citations: the retained excerpts carry no citation command at all, so no bibliography entry of the article is transcribed and no reference is redistributed. Reference adaptation: none was needed and none was made; every reference inside the delivered unit resolves inside it, so no unresolved reference or citation is produced. Printed slips kept literal: no printed word, symbol, hypothesis or number of the counted statement or of its proof was altered. Layout and class: the article's own class is \texttt{amsart} with packages including geometry, lipsum, amsfonts, graphicx, epstopdf, pifont, enumitem, tikz, pgfplots, pgfplotstable, tkz-base, tkz-euclide, hyperref, xcolor; the wrapper is the lane's pinned \texttt{amsart} editorial wrapper at 10pt on A4, which restates the theorem-like environments the retained text needs and adds the article's own macro definitions that the retained text needs (\texttt{\textbackslash Bern}, \texttt{\textbackslash Fluxnum}, \texttt{\textbackslash H}, \texttt{\textbackslash R}, \texttt{\textbackslash dt}, \texttt{\textbackslash h}, \texttt{\textbackslash v}), with extra wrapper packages (bm, bbm, dsfont, xspace, cancel, mathtools, enumitem). The delivered numbering is the unit's own. Assets: no retained span contains a figure environment, a graphics-inclusion command, a PDF-page inclusion, a table or a TikZ picture; no image file is copied into this package, the package declares no asset dependency and no image file is redistributed with it. Licence: version arXiv:2509.14908v1 of this article is distributed under the Creative Commons Attribution 4.0 International licence, as recorded in the arXiv licence metadata for this exact version and shown on the abstract page https://arxiv.org/abs/2509.14908v1. A full rights notice is delivered at licenses/arxiv-2509.14908-CC-BY-4.0.txt. Characters: every retained excerpt is pure ASCII, so the delivered engine, XeLaTeX with its default Latin Modern text font, reports no missing character.
\begin{subequations} \label{eq:scheme}
\begin{align}
	\label{eq:scheme_u}
		h_i \frac{L^n u_i^n - L^{n-1} u_i^{n-1}}{\dt} + \Fluxnum{i+} - \Fluxnum{i-} &= 0, \quad \textnormal{ for } 1 \leq i \leq I, \\
	\label{eq:left_cond_flux}
		\Fluxnum{} - a + bu_0^n &= 0, \\
	\label{eq:right_cond_flux}
		\Fluxnum{I+} &= 0, \\
	\label{eq:X0}
	\frac{X_0^n - X_0^{n-1}}{\dt} - \alpha_0 + \beta_0 u_0^{n} - (1-R) \frac{X_1^n - X_1^{n-1}}{\dt} &= 0, \\
	\label{eq:X1}
		\frac{X_1^n - X_1^{n-1}}{\dt} + \alpha_1 - \beta_1 u_{I+1}^{n} &= 0, \\
		L^n - X_1^n + X_0^n &= 0.\label{eq:L}
\end{align}
\end{subequations}

	\begin{equation}
	    \label{eq:def_v}
		\v{i+}{n} = (1-R) \delta^n [X_1] - \xi_{i+\frac{1}{2}} \delta^n [L] - \delta^n [X_0].
	\end{equation} 

	\begin{lemma} \label{lem:Bern}
		The Bernoulli function is $1$-Lipschitz continuous, positive, decreasing and satisfies:
		\begin{enumerate}[label=(\textit{\roman*})]
			\item For $r \in \R$, $B(-r) - B(r) = r$;
			\item For $r \in \R$, $B(r)e^r = B(-r)$;
			\item For $\theta \in [0,1]$, $p,q,r \in \R$:
				\begin{equation} \label{eq:Bern}
					B(-r)p - B(r)q = (\theta B(r) + (1-\theta) B(-r))(p-q) + r((1-\theta)q + \theta p).
				\end{equation}
		\end{enumerate}
	\end{lemma}

	\begin{remark} \label{rk:pi_croiss}
		Since $\pi_\phi'(r) = r \phi''(r)$ with $\phi$ convex, $\pi_\phi$ is non-decreasing over $\R_{\geq 0}$.
	\end{remark}

	\begin{definition}
		Let $\phi \in \mathcal{C}^2(\R)$ an arbitrary convex function. Set $n \geq 0$. The discrete free energy associated to $\phi$ is given by
		\begin{equation*}
			\H^n = \sum_{i=1}^I L^n h_i \phi(u_i^n).
		\end{equation*}
		Moreover, the discrete total free energy associated to $\phi$  is
		\begin{multline} \label{eq:def_htot_disc}
			\H^{{\text{tot}},n} = \H^n + \pi_\phi \left( \frac{\alpha_0}{\beta_0} \right) \sum_{k=1}^n \dt (\alpha_0 - \beta_0 u_0^k) \\
			+ \phi' \left(\frac{a}{b} \right) \sum_{k=1}^n \dt (a - b u_0^k) + R\pi_\phi \left( \frac{\alpha_1}{\beta_1} \right) \sum_{k=1}^n \dt (\alpha_1 - \beta_1 u_{I+1}^k).
		\end{multline}
	\end{definition}

	\begin{proposition} \label{prop:disc_energ_dec}
	    Assume that the scheme \eqref{eq:scheme} admits a solution $((u_i^n)_{0 \leq i \leq I+1},$ $X_0^n, X_1^n)_{0 \leq n \leq N} $ such that, for all $0\leq n\leq N$, $L^n>0$ and $u_i^n\geq 0 $ for all $0\leq i\leq I+1$.
		Then, for any convex function $\phi \in \mathcal{C}^2(\R)$, there exists a non-negative dissipation term $\mathcal{D}_{\phi}^n$ such that the discrete derivative of $\H^{{\text{tot}},n}$ satisfies
		\begin{equation} \label{eq:disc_htot_decr}
			\dfrac{\H^{{\text{tot}},n} - \H^{{\text{tot}},n-1}}{\dt} + \mathcal{D}_{\phi}^n \leq 0.
		\end{equation}
	\end{proposition}

	\begin{proof}
	Let us introduce the application $A: (l,w) \mapsto l \phi\left(\dfrac{w}{l}\right)$ in order to rewrite 
	$$
	\H^n=\sum_{i=1}^I h_i A\left(L^n, L^n u_i^n \right).
	$$
	As $\phi$ is convex, $A$ is jointly convex and verifies $\nabla A(l,w)=(-\pi_\phi(\dfrac{w}{l}), \phi'(\dfrac{w}{l}))$. 
	From an elementary convexity inequality, we deduce that 
		\begin{equation}\label{disc_energ_dissip_startpt}
			\frac{\H^n - \H^{n-1}}{\dt} \leq Y^n + Z^n
		\end{equation}
		with
		\[
			Y^n = \sum_{i=1}^I h_i \phi'(u_i^n) \frac{L^n u_i^n - L^{n-1} u_i^{n-1}}{\dt}, \qquad 
			Z^n =  -\delta^n[L] \sum_{i=1}^I h_i  \pi_\phi(u_i^n).
		\]
		Using the numerical scheme \eqref{eq:scheme_u} in the expression $Y^n$, performing a discrete integration by parts and using the 
		discrete boundary conditions~\eqref{eq:left_cond_flux} and \eqref{eq:right_cond_flux} provides 
		\begin{equation}\label{def:Yn}
			Y^n  = (a-b u_0^n) \phi'(u_0^n) + \sum_{i=0}^{I} \left( \phi'(u_{i+1}^n) - \phi'(u_{i}^n) \right) \Fluxnum{i+}.
		\end{equation}
		For any $0 \leq i \leq I$ and $n \geq 0$, the mean value formula and the link between $\phi$ and $\pi_\phi$ 
		ensure the existence of $\theta_{i+\frac{1}{2}}^n \in [0,1]$ such that
		\begin{equation} \label{eq:mean_value}
			\pi_\phi(u_{i+1}^n) - \pi_\phi(u_{i}^n) = \left(\theta_{i+\frac{1}{2}}^n u_{i}^n + (1-\theta_{i+\frac{1}{2}}^n)u_{i+1}^n\right) \left(\phi'(u_{i+1}^n) - \phi'(u_{i}^n)\right).
		\end{equation}
	Then, using the third property of the Bernoulli function given in \eqref{eq:Bern} in Lemma \ref{lem:Bern}, it comes that
		\begin{equation} \label{eq:Yn-2}
			\sum_{i=0}^I\left( \phi'(u_{i+1}^n) - \phi'(u_{i}^n) \right) \Fluxnum{i+} =
			- \mathcal{D}_{\phi}^{\text{bulk},n} + \sum_{i=0}^I v_{i + \frac{1}{2}}^n (\pi_\phi(u_{i+1}^n) - \pi_\phi(u_{i}^n)),
		\end{equation}
		where  the quantity $\mathcal{D}_{\phi}^{\text{bulk},n}$ is defined by
		\begin{multline}\label{eq:dbulk.sum}
		\mathcal{D}_{\phi}^{\text{bulk},n} = \sum_{i=0}^{I}
		\left(\Bern{i+}{n}{} \theta_{i+\frac{1}{2}}^n + \Bern{i+}{n}{-}(1 - \theta_{i+\frac{1}{2}}^n) \right)\times  \\
			\frac{\phi'(u_{i+1}^n) - \phi'(u_{i}^n) }{L^n \h{i+}} (u_{i+1}^n - u_{i}^n)  \geq 0.
	\end{multline}
	Thanks to a discrete integration by parts, the last term in \eqref{eq:Yn-2} rewrites 
	$$
	\sum_{i=0}^I v_{i + \frac{1}{2}}^n (\pi_\phi(u_{i+1}^n) - \pi_\phi(u_{i}^n))=-\sum_{i=1}^I (v_{i + \frac{1}{2}}^n-v_{i - \frac{1}{2}}^n)\pi_\phi (u_{i}^n) + v_{I+\frac12}^n \pi_\phi(u_{I+1}^n) -v_{\frac{1}{2}}^n\pi_\phi(u_0^n).
	$$
	The definition \eqref{eq:def_v} of the discrete velocities implies that \((v_{i + \frac{1}{2}}^n-v_{i - \frac{1}{2}}^n)=-\delta^n[L]h_i\) and also fix values of  $v_{I+\frac{1}{2}}^n$ and $v_{\frac{1}{2}}^n$,
	so that 
	\begin{multline}\label{est_termesadd}
	\sum_{i=0}^I v_{i + \frac{1}{2}}^n (\pi_\phi(u_{i+1}^n) - \pi_\phi(u_{i}^n))=-Z^n-R\delta^n[X_1]\pi_\phi (u_{I+1}^n)\\
	-((1-R)\delta^n[X_1]-\delta^n[X_0])\pi_\phi(u_0^n).
	\end{multline}
		Then, using the equations on $\delta^n[X_0]$ and $\delta^n[X_1]$ in \eqref{eq:scheme}, we deduce from \eqref{disc_energ_dissip_startpt}, \eqref{def:Yn}, \eqref{eq:Yn-2} and \eqref{est_termesadd}:
		\begin{multline*}
			\dfrac{\H^n - \H^{n-1}}{\dt} + \mathcal{D}_{\phi}^{\text{bulk},n} \\
			+ (bu_0^n - a)\phi'(u_0^n)+ (\beta_0 u_0^n-\alpha_0 )\pi_\phi(u_0^n) 
			 + R( \beta_1 u_{I+1}^n-\alpha_1)\pi_\phi(u_{I+1}^n) \leq 0.
		\end{multline*}
		
		In order to recover the expression of $\H^{{\text{tot}},n}$ given in \eqref{eq:def_htot_disc}, it only remains to use the non-decreasing behavior of $\phi'$ and $\pi_\phi$ (see Remark \ref{rk:pi_croiss}). Then, the following term is non-negative:
		\begin{multline}\label{eq:dbound}
			\mathcal{D}^{\text{bound}, n}_\phi = (\beta_0 u_0^n-\alpha_0 )\left(\pi_\phi(u_0^n) - \pi_\phi\left(\frac{\alpha_0}{\beta_0}\right)\right) \\
			+ (bu_0^n-a)\left(\phi'(u_0^n) - \phi'\left(\frac{a}{b}\right)\right) \\
			+R(\beta_1 u_{I+1}^n-\alpha_1)\left(\pi_\phi(u_{I+1}^n) - \pi_\phi\left(\frac{\alpha_1}{\beta_1}\right)\right) \geq 0.
		\end{multline}
		Finally, the  inequality \eqref{eq:disc_htot_decr} holds
		with the dissipation term $\mathcal{D}_{\phi}^n$ defined as 
		\begin{equation} \label{eq:disc_diss_term}
			\mathcal{D}_{\phi}^n =  \mathcal{D}_{\phi}^{\text{bulk},n} + \mathcal{D}^{\text{bound}, n}_\phi.
		\end{equation}
	\end{proof}

\end{document}
